\documentclass[addpoints]{exam}
% \documentclass[addpoints,12pt]{exam}

\usepackage[slovene]{babel}
\usepackage{amsfonts}
\usepackage{amssymb}
\usepackage{amsmath}
\usepackage{float}
\usepackage{graphicx}
\usepackage{enumitem}
\usepackage{color}
\usepackage[gen]{eurosym}
\usepackage{newunicodechar}
\newunicodechar{€}{\euro}

%Komentar naslednje vrstice glede na verzijo testa -- rešitve ali prostor za odgovore:
\printanswers

% \linespread{2}


\pointsinrightmargin
\marginbonuspointname{}


\renewcommand{\solutiontitle}{\noindent\textbf{Rešitve in točkovanje:}\par\noindent}

\colorgrids
\definecolor{GridColor}{gray}{0.7}
\setlength{\gridsize}{7mm}

\extrawidth{5mm}
\setlength{\rightpointsmargin}{1.5cm}

\begin{document}

\runningfooter{}{}{}

\begin{center}
    \fbox{\fbox{\parbox{15cm}{\centering
    Tretje pisno ocenjevanje znanja \\ 1. letnik -- test A \\ 9. februar 2026 \\ (Racionalna števila)}}}
\end{center}

\vspace{0.1in}
\makebox[\textwidth]{Ime in priimek, oddelek:\enspace\hrulefill}


\begin{center}
    \chqword{Naloga}
    \chpword{Možne točke}
    \chbpword{Dodatne točke}
    \chsword{Dosežene točke}
    \chtword{Skupaj}  
    \combinedpointtable[h][questions]
    % \combinedgradetable[h][questions]

\end{center}

\vspace{0.1in}


\begin{questions}

    % 1. vprašanje

    \question
    Izračunajte natančno vrednost izraza.
    \begin{parts}
        \part[4]
        $\dfrac{\frac{2}{3}+1\frac{5}{6}\cdot 3}{\frac{2}{7}:\frac{4}{3}-1\frac{1}{14}+2}$

        \begin{solution}[8cm]
            \begin{align*} 
                \dfrac{\frac{2}{3}+1\frac{5}{6}\cdot 3}{\frac{2}{7}:\frac{4}{3}-1\frac{1}{14}+2}&=\dfrac{\frac{2}{3}+\frac{11}{6}\cdot 3}{\frac{2}{7}\cdot\frac{3}{4}-\frac{15}{14}+\frac{28}{14}}= \\
                                                                                                &=\dfrac{\frac{2}{3}+\frac{11}{2}}{\frac{1}{7}\cdot\frac{3}{2}+\frac{13}{14}}= \\
                                                                                                &=\dfrac{\frac{4}{6}+\frac{33}{6}}{\frac{3}{14}+\frac{13}{14}}= \\
                                                                                                &=\dfrac{\frac{37}{6}}{\frac{16}{14}}= \\
                                                                                                &=\dfrac{37\cdot 14}{6\cdot 16}= \\
                                                                                                &=\dfrac{259}{48} 
            \end{align*}
            Za pravilno seštevanje/odštevanje ulomkov $1$ točka; pravilno množenje ulomkov $1$ točka; pravilno razrešen dvojni ulomek $1$ točka, okrajšan rezultat $1$ točka.
        \end{solution}


        \part[5]
        $0.1\overline{3333}:\left(1.2-0.6\cdot\left(0.04+0.06\right)\right)^{-1}$

        \begin{solution}[4cm]
            \begin{align*}
                0.1\overline{3333}:\left(1.2-0.6\cdot\left(0.04+0.06\right)\right)^{-1} &= \dfrac{2}{15}\cdot \left(1.2-0.6\cdot 0.1\right)= \\
                                                                                        &= \dfrac{2}{15}\cdot \left(1.2-0.06\right)= \\
                                                                                        &= \dfrac{2}{15}\cdot 1.14= \\
                                                                                        &= \dfrac{2}{15}\cdot \dfrac{114}{100}= \\
                                                                                        &= \dfrac{19}{125}
            \end{align*}
            Za pravilen zapis periodičnega števila z ulomkom $1$ točka, pravilno seštevanje/odštevanje $1$ točka, pravilno množenje $1$ točka, upoštevanje deljenja in obratne vrednosti $1$ točka, okrajšan rezultat $1$ točka.
        \end{solution}

    \end{parts}






    % \newpage

        % 3. vprašanje

        \question
        Poenostavite izraze.
        \begin{parts}
            

            \part[7]
            $\left(\dfrac{6-x}{x^2+6x}-\dfrac{x}{36-x^2}\right)\left(\dfrac{2x-6}{x^2+6x}\right)^{-1}+\dfrac{x}{6-x}$
    
            \begin{solution}[9cm]
                \begin{align*}
                    \left(\dfrac{6-x}{x^2+6x}-\dfrac{x}{36-x^2}\right)\left(\dfrac{2x-6}{x^2+6x}\right)^{-1}+\dfrac{x}{6-x} &= \left(\dfrac{6-x}{x(x+6)}-\dfrac{x}{(6-x)(6+x)}\right)\dfrac{x(x+6)}{2(x-3)}+\dfrac{x}{6-x}= \\
                                                                                                                            &= \left(\dfrac{(6-x)^2}{x(x+6)(6-x)}-\dfrac{x^2}{x(6-x)(6+x)}\right)\dfrac{x(x+6)}{2(x-3)}+\dfrac{x}{6-x}= \\
                                                                                                                            &= \dfrac{36-12x+x^2-x^2}{x(x+6)(6-x)}\dfrac{x(x+6)}{2(x-3)}+\dfrac{x}{6-x}= \\
                                                                                                                            &= \dfrac{12(3-x)}{x(x+6)(6-x)}\dfrac{x(x+6)}{2(x-3)}+\dfrac{x}{6-x}= \\
                                                                                                                            &= \dfrac{-6}{6-x}+\dfrac{x}{6-x}= \\
                                                                                                                            &= \dfrac{-6+x}{6-x}= \\
                                                                                                                            &= -1
                \end{align*}
                Za pravilno faktorizacijo izrazov $1$ točka, pravilno razčlenjevanje izrazov $1$ točka, pravilno seštevanje ulomkov $1$ točka, pravilna obratna vrednost ulomka $1$ točka, pravilno množenje ulomkov $1$ točka, pravilno krajšanje ulomkov $1$ točka, okrajšan rezultat $1$ točka.
            \end{solution}
    
    
            \part[5]
            $\left(\dfrac{x^{-3}-x^{-1}}{1-x^{-2}}\right)^{-1}+\left(\dfrac{1}{x}\right)^{-1}$
    
            \begin{solution}[7cm]
                \begin{align*}
                    \left(\dfrac{x^{-3}-x^{-1}}{1-x^{-2}}\right)^{-1}+\left(\dfrac{1}{x}\right)^{-1} &= \left(\dfrac{\frac{1}{x^3}-\frac{1}{x}}{1-\frac{1}{x^2}}\right)^{-1}+x= \\
                                                                                                    &= \left(\dfrac{\frac{1}{x^3}-\frac{x^2}{x^3}}{\frac{x^2}{x^2}-\frac{1}{x^2}}\right)^{-1}+x= \\
                                                                                                    &= \left(\dfrac{\frac{1-x^2}{x^3}}{\frac{x^2-1}{x^2}}\right)^{-1}+x= \\
                                                                                                    &= \left(\dfrac{(1-x^2)x^2}{(x^2-1)x^3}\right)^{-1}+x= \\
                                                                                                    &= \left(-\dfrac{1}{x}\right)^{-1}+x= \\
                                                                                                    &= -x+x= \\
                                                                                                    &= 0
                \end{align*}
                Za pravilen zapis obratnih vrednosti ulomkov $1$ točka,  pravilno seštevanje ulomkov $1$ točka, pravilno razrešen dvojni ulomek $1$ točka, pravilno krajšanje ulomkov $1$ točka, končen rezultat $1$ točka.
            \end{solution}
    
    
            \part[5]
            $\left(\dfrac{-3^4x^{-2}y^3}{x^3z^2}\right)^{-4} : \left(\dfrac{y^{-3}}{3^5x^2z^{-2}}\right)^3$    
            \begin{solution}[2cm]
                \begin{align*}
                    \left(\dfrac{-3^4x^{-2}y^3}{x^3z^2}\right)^{-4} : \left(\dfrac{y^{-3}}{3^5x^2z^{-2}}\right)^3 &= \dfrac{3^{-16}x^{8}y^{-12}}{x^{-12}z^{-8}} : \dfrac{y^{-9}}{3^{15}x^6z^{-6}}= \\
                                                                                                                    &= \dfrac{3^{-16}x^{8}y^{-12}}{x^{-12}z^{-8}} \cdot \dfrac{3^{15}x^6z^{-6}}{y^{-9}}= \\
                                                                                                                    &= 3^{-16+15}x^{8-(-12)+6}y^{-12-(-9)}z^{-(-8)+(-6)}= \\
                                                                                                                    &=\frac{x^{26}z^2}{3y^3}
                \end{align*}
                Za pravilno potenciranje z negativnim eksponentom $1$ točka, pravilno potenciranje s pozitivnim eksponentom $1$ točka, pravilno deljenje $1$ točka, pravilno krajšanje $1$ točka, končen okrajšan rezultat $1$ točka.
            \end{solution}
    
    
        \end{parts}
   
   
   
   
        % \newpage
    % 2. vprašanje

    \question
    Okrajšajte ulomke.
    \begin{parts}
        

        \part[3]
        $\dfrac{2^{n-1}+3\cdot 2^n}{4^n+5\cdot 2^{2n-1}}$

        \begin{solution}[5cm]
            \begin{align*}
                \dfrac{2^{n-1}+3\cdot 2^n}{4^n+5\cdot 2^{2n-1}} &= \dfrac{2^{n-1}(1+3\cdot 2)}{2^{2n-1}(2+5)}= \\
                                                                &= \dfrac{2^{n-1}}{2^{2n-1}}= \\
                                                                &= 2^{n-1-(2n-1)}= \\
                                                                &= \dfrac{1}{2^n}
            \end{align*}
            Pravilno izpostavljanje $1$ točka, pravilno krajšanje $1$ točka, končen okrajšan rezultat $1$ točka.
        \end{solution}


        \part[4]
        $\left(\dfrac{\dfrac{2}{x^{-2}}}{\left(x^{-2}\right)^{-1}}\right)^{-1}$

        \begin{solution}[6cm]
            \begin{align*}
                \left(\dfrac{\dfrac{2}{x^{-2}}}{\left(x^{-2}\right)^{-1}}\right)^{-1} &= \dfrac{\left(x^{-2}\right)^{-1}}{\dfrac{2}{x^{-2}}}= \\
                                                                                    &= \dfrac{x^2}{\dfrac{2}{x^{-2}}}= \\
                                                                                    &= \dfrac{x^2\cdot x^{-2}}{2}= \\
                                                                                    &= \dfrac{1}{2}
            \end{align*}
            Za pravilno zapisano obratno vrednost $1$ točka, pravilno potenciranje z negativnim eksponentom $1$ točka, pravilno razreševanje dvojnega ulomka $1$ točka, končen okrajšan rezultat $1$ točka.
        \end{solution}



    \end{parts}



        \question[4]
        Kolikšno koncentracijo raztopine dobimo, če zmešamo $2~kg$ $60~\%$ raztopine in $3~kg$ $40~\%$ raztopine?

        \begin{solution}[3cm]
            \begin{align*}
                0.60 \cdot 2~kg + 0.40\cdot 3~kg &= x \cdot 5~kg \\
                                 1.2~kg+1.2~kg &= x\cdot 5~kg \\
                                            x &= \dfrac{2.4~kg}{5~kg} \\
                                            x &= 0.48
            \end{align*}
            Pravilna nastavitev enačbe $1$ točka, pravilno reševanje $1$ točka, pravilen rezultat $1$ točka, zapis pravilnega odgovora $1$ točka.
        \end{solution}



        % \newpage


        \question[3]
        Servis so najprej podražili za $10~\%$, potem pa se je ena skodelica okrušila in so
        ga pocenili za $30~\%$. Koliko je servis stal na začetku, če ga danes lahko kupimo za $115.5~\euro$?

        \begin{solution}[6cm]
            \begin{align*}
                0.70\cdot \left(1.10\cdot x\right) &=115.5~\euro \\
                                     0.77 \cdot x &= 115.5~\euro \\
                                                x &= \dfrac{115.5}{0.77}~\euro \\
                                                x &= 150~\euro
            \end{align*}
            Pravilen zapis in reševanje enačbe $1$ točka, pravilen rezultat $1$ točka, zapis pravilnega odgovora $1$ točka.
        \end{solution}
     ~\\


     % dodatno vprašanje
    \bonusquestion[4]
    \textit{Dodatna naloga:} Poenostavite izraz. $$\left(1+\dfrac{2}{x+\dfrac{2}{2+\frac{3}{x+1}}}\right):\dfrac{1}{\dfrac{2x^2+7x+2}{x^3+64}}$$

    \begin{solution}[7cm]
        \begin{align*}
            \left(1+\dfrac{2}{x+\dfrac{2}{2+\frac{3}{x+1}}}\right):\dfrac{1}{\dfrac{2x^2+7x+2}{x^3+64}} &= \left(1+\dfrac{2}{x+\dfrac{2}{\frac{2x+2+3}{x+1}}}\right)\cdot \dfrac{2x^2+7x+2}{x^3+64}= \\
                                                                                                        &= \left(1+\dfrac{2}{x+\dfrac{2x+2}{2x+5}}\right)\cdot \dfrac{2x^2+7x+2}{x^3+64}= \\
                                                                                                        &= \left(1+\dfrac{2}{\dfrac{2x^2+5x+2x+2}{2x+5}}\right)\cdot \dfrac{2x^2+7x+2}{x^3+64}= \\
                                                                                                        &= \left(1+\dfrac{4x+10}{2x^2+7x+2}\right)\cdot \dfrac{2x^2+7x+2}{x^3+64}= \\
                                                                                                        &= \dfrac{2x^2+7x+2+4x+10}{2x^2+7x+2}\cdot \dfrac{2x^2+7x+2}{x^3+64}= \\
                                                                                                        &= \dfrac{2x^2+11x+12}{x^3+64}= \\
                                                                                                        &= \dfrac{(2x+3)(x+4)}{(x+4)(x^2-4x+16)}= \\
                                                                                                        &= \dfrac{2x+3}{x^2-4x+16}= \\
        \end{align*}
        Za pravilno reševanje izraza do $3$ točke; končen rezultat $1$ točka.
    \end{solution}




\end{questions}



\end{document}